% !TeX root = ../document_maitre.tex

Patrimoine et gestion des collections sont des domaines étroitement liés, actuellement en proie à des difficultés exarcerbée et, pense-t-on, à l'aube d'une révolution technologique. 

La double hétérogénéité des collections, des données informatiques générées ainsi que leurs masses sont aujourd'hui des problématiques fondamentales. Chaque typologie patrimoniale connaît ses propres normes et standards de descriptions, ainsi que ces propres modèles conceptuel de gestions. Leurs numérisation dupplique les obligations et besoins de gestion, comme nous l'avons vus avec l'\gls{imgani} où cette double nature amène de doubles enjeux de conservations. Et enfin, leurs masses, dûs au phénomène de \gls{bgdt}, complexifie les efforts d'harmonisation et de gestions. 

Les \gls{sgc} sont au \coeur de nos pratiques de gestion des collections. Plus que de simple registre numérique, ce sont de véritables outils stratégiques afin de piloter la vie des collections et institutions conservatrices. L'évolution des difficultés et des besoins dans ce domaine ont toujours amenée ces logiciels à s'adapter. L'intégration des normes et standards, l'applications de l'\gls{intero} et des modèles conceptuels sont aujourd'hui prise en compte, comme nous l'avons vus avec la suite SkinSoft.

Mais aujourd'hui, la \gls{bgdt} amène un renouveau des difficultés et ne ces dernier ne semblent plus disposer des moyens de s'y adapter pleinement. Le problème majeur se situe en amont du logiciel : lors de la production et de l'intégration des données. Par exemple, pour l'\gls{imgani} les \gls{sgc} ne proposent pas de solution concrète de facilitation d'une indexation fastidieuse. Situé au \coeur de ces derniers, l'indexation humaine de telles typologies patrimoniales est ajourd'hui limitée. Elle ne permet plus de traiter à un rythme suffisant la masse de données entrante. Ajoutons que les modalités de gestion très spécifiques et maléables de ces typologies rendent leurs intégrations de manière simple et ergonomique assez complexe à établir. 

Ces logiciels ont alors deux limites principales face aux évolutions des pratiques aux enjeux du \gls{bgdt} dans notre milieu. D'abord fonctionnelle, car ils n'arrivent plus à faciliter les processus d'indexation et de gestion des données entrantes dans des modèles de gestion complexes. Ensuite ergonomiques, car la nature complexe de ces logiciels n'offre plus de moyen de repères suffisant afin d'aborder sereinement cet outils de travail.\newline

Aujourd'hui, l'\gls{IA} est en plein essort et se fait remarquer par sa capacité de traitement en autonomie de très grandes volumétries de données. Mais également pour sa capacité d'interraction avec d'autres logiciels. 

Dans le domaine patrimonial, nous solicitons surtout les technologies de \gls{NLP}, \gls{RAG} et de \gls{CV}. L'utilisation de ces outils est paradoxale : rapidement utilisé de différentes manière pour le patrimoine, la recherche semble aujourd'hui préférer une utilisation pour l'accessibilité grand public. Leurs applications à la sphère professionnelles pour la production de métadonnées s'oppose à des principes de rigueur scientifique. En effet, leurs capacités à halluciner constitue un risque de corruption du bien fondé des bases de données institutionnelles. Vraisemblablement, cela amène à repousser le moment d'intégration de ces outils.

Cependant, ces technologies propose des pistes prometteuses pour soulager, si ce n'est résoudre, la plus parts des difficultés actuelles des \glspl{sgc}. Le \gls{RAG} et le \gls{NLP} sont prometteur afin de permettre des recherches plus naturelles, fines et faciles dans sa propre base de données. La capacité de ce dernier à \guillemet{comprendre} le langage humain semble en faire un parfait intermédiare entre l'application et l'humain. Pour la \gls{CV}, c'est sa capacité à analyser des images, identifier et exploiter des caractéristiques clés dans des résumés sémantiques qui représente une voie d'amélioration de l'indexation de l'\gls{imgani}. 

Ces outils peuvent répondre aux difficultés fonctionnelles, par des processus d'indexation automatisée en interaction avec des \gls{thesaurus} par exemple. Mais aussi aux difficultés ergonomiques en facilitant l'accessibilité à l'outil et aux données, peut être même jusqu'à appuyer l'homogénéisation des données. De tels outils semblent être prometteurs pour appuyer les professionnels dans leurs diverses tâches de production et d'utilisation des \glspl{sgc}. \newline

Ces outils n'ont pas encore été pensés pour une intégration au \coeur de \gls{sgc}. Cela pose donc la question de la faisabilité technique et de l'impact potentiel sur les pratiques et fonctionnement de ces logiciels. 

Nous avons donc suivi cette piste en se demandant comment intégrer de tels technologies dans le logiciel SkinSoft, dans quel but et surtout quels en sont les enjeux. Cela à aboutit à explorer les deux axes à l'aide d'outils pensé pour appuyer les professionnels. D'abord l'axe ergonomique avec la \guillemet{recherche augmentée} pour une amélioration de la pertinence des résultats et une facilité d'interrogation des données institutionnelles. Mais également avec le \gls{chatbot}, pensé pour accompagner l'utilisateur dans son utilisation de l'application et ses missions de contrôle de la qualité des métadonnées indexées. Mais également l'axe fonctionnel avec l'indexation semi-automatisées des \glspl{imgani}, représentant un volume très importants de données à intégrer selon des exigences normatives variables et exigeante. 

Nos différentes explorations, encore très conceptuelles, de ces outils nous amène à différents constats. D'abord, l'inégalité des contraintes techniques et des faisabilités technologiques. Des processus comme la \guillemet{recherche augmentée} sont aujourd'hui basés sur des technologies plus stables et moins gourmandes en ressources qu'un processus d'indexation semi-automatisée de l'\gls{imgani}. Ce dernier semblant dès lors bien plus inaccessibles que la \guillemet{recherche augmentée}, à cause de sa complexité et de sa demande en ressources complexe à satisfaire. 

Nous démontrons que la réussite de ces outils ne dépendent pas de ces seuls question. Ce sont en effet des conditions préalables à la réalisations de ces outils, mais l'enjeu final de l'adoption de l'outil réside dans d'autre domaine. À savoir l'\gls{ux} et l'\gls{ui}, en charge de la facilité d'utilisation de ces outils. De cette facilité d'usage dépend l'adoption de ces outils par les professionnels du patrimoine. C'est finalement ce degré d'adoption de l'outil qui transforme ou non nos pratiques. Ce qui finalement place les questions d'\gls{ux} et d'\gls{ui} au \coeur des enjeux. 

De plus, de tels outils s'inscrivent dans un cadre réglementaire et technique strict. Le patrimoine et la gestion des données sont riches en contraintes et exigence. La souveraineté des informations, la protection du droit d'auteur, et la conformité au cadre juridique de l'\textit{AI Act} sont des enjeux capitaux. Plus que la performance, un bon outil d'appui aux professionnels dans un \gls{sgc} est avant tout un outil respecteux et conforme à cet univers.

Sur le plan éthique, l'intégration de tels outils est aussi problématique. Ils questionnent des valeurs d'honnêteté, de rigueur et d'exhaustivité scientifique dans les bases de données institutionnelles. Combinés aux facteurs juridiques, cela réaffirme que l'automatisation ne peut être totale. L'approche de ces outils doit être évolutive et prudente. Cela passe par des architectures \textit{(human-in-the-loop)}, demeurant la condition \textit{sine qua non} d'une exploitation éthique, pérenne et acceptée.\newline

Finalement, l'implémentation d'outils d'\gls{IA} peu permettre de répondre aux défis actuels de deux manières. D'abord par l'augmentation de l'accessibilité et de la facilité d'usage des \gls{sgc} et des bases de données patrimoniales. Bien qu'assez largement ignorés aujourd'hui, nous avons démontré que la première limite de ces logiciels viens de leurs propres architecture logicielle. Les \glspl{chatbot} et recherche assistées par \gls{IA} peuvent permettre d'en améliorer l'accessibilité, d'en facilité l'usage et d'offrir des expériences d'utilisations plus confortables créant de meilleur conditions de productions. 

Ensuite par l'automatisation relative de chaîne d'indexation sur certaines typologies documentaires. Bien que notre essais sur l'\gls{imgani} ai montré toute la complexité et les limites d'un tel processus, on peut penser qu'appliqué à des objets plus structuré et traités comme des archives, de tels processus puisse être réalisable. Cependant, la consommation de ressources de ces outils et les danger de l'hallucination demandent à en relativiser l'application. \newline

Cette étude nous permet finalement de répondre à notre problématique initiale. Ces technologies peuvent améliorer nos utilisations des \glspl{sgc} et nos capacités d'indexations, sans pour autant les révolutionner étant donné leurs limites. Ainsi, ils ne peuvent remplacer l'expertise humaine. L'arrivée de ces outils dans tels domaines de production amènent des réflexions éthiques et juridique qui convergent vers leurs rationnalisation pour en faire des outils d'appuis, de facilitation et non de remplacement. 

Loin de \guillemet{dénaturer} nos pratiques, il est naturel que ces dernières évoluent sans pour autant changer drastiquement. Et il est aujourd'hui de notre responsabilité d'ériger ces cadres éthiques, juridiques et technologiques pour une conception et une utilisation raisonnée de tels outils, cruciaux pour répondre aux défis de la gestion des collections.
